\chapter{Documentación Legal y Administrativa}

\section{Licencias y autorizaciones}

Para garantizar el inicio oportuno de las actividades y evitar paralizaciones por incumplimiento normativo, se gestionarán y presentarán los siguientes permisos antes y durante la ejecución del proyecto:

\subsection*{Autorizaciones Institucionales (UNSAAC)}

Dado que el proyecto se desarrolla dentro del Campus Universitario de Perayoc, se tramitarán los siguientes permisos internos ante la Dirección de la Escuela Profesional y la Oficina de Infraestructura:

\begin{enumerate}
    \item \textbf{Permiso de Acceso y Permanencia:} Autorización nominal para el ingreso del personal técnico, herramientas y materiales al pabellón, incluyendo permisos especiales para trabajos en horarios extendidos o fines de semana si el cronograma lo requiere.
    \item \textbf{Autorización de Intervención Física:} Permiso emitido por la Unidad de Mantenimiento/Infraestructura para realizar perforaciones, fijación de canaletas y uso de los ductos montantes existentes en los \totalPisos{} niveles del edificio, incluyendo el paso de los cables del Nivel 1 hacia el MDF del Nivel 2.
    \item \textbf{Coordinación con T.I.:} Acta de conformidad con el Centro de Cómputo o responsable de TI de la Facultad para acceder a los cuartos de telecomunicaciones (IDF/MDF) existentes y realizar las conexiones al backbone de fibra óptica.
\end{enumerate}

\subsection*{Licencias de Seguridad y Laborales (Normativa Peruana)}

En cumplimiento de la Ley N.º 29783 (Ley de Seguridad y Salud en el Trabajo), es requisito indispensable para el inicio de obra presentar:

\begin{enumerate}
    \item \textbf{SCTR (Seguro Complementario de Trabajo de Riesgo):} Pólizas vigentes de Salud y Pensión para todo el personal técnico y operativo que ingresará a la obra, con cobertura para trabajos de riesgo (electricidad y altura).
    \item \textbf{Registro de Inducción y Entrega de EPP:} Documentación firmada que acredite que los trabajadores han recibido la charla de inducción de seguridad específica para el sitio y cuentan con los Equipos de Protección Personal normados.
    \item \textbf{IPERC y ATS:} Matriz de Identificación de Peligros y Evaluación de Riesgos (IPERC) y los Análisis de Trabajo Seguro (ATS) diarios, firmados por el prevencionista de riesgos.
\end{enumerate}

\subsection*{Consideraciones Municipales y Ambientales}

\begin{itemize}
    \item \textbf{Licencia de Obra:} De acuerdo con los alcances del proyecto definidos en la Memoria Descriptiva, los trabajos se clasifican como ``Acondicionamiento y Refacción'' de instalaciones interiores. Al no modificar la estructura, fachada, ni el volumen del edificio, no se requiere tramitar una Licencia de Edificación Mayor ante la Municipalidad, bastando con las autorizaciones internas de la Universidad.
    \item \textbf{Gestión de Residuos:} Se contará con la autorización para el retiro y disposición final de residuos de construcción (escombros de perforación) y residuos de aparatos eléctricos y electrónicos (RAEE) -como cables antiguos retirados- a través de una Empresa Operadora de Residuos Sólidos (EO-RS) autorizada, cumpliendo con la normativa ambiental vigente.
\end{itemize}

\section{Permisos de uso de infraestructura}

La ejecución del proyecto requiere la ocupación temporal y permanente de espacios físicos dentro del Pabellón de Ingeniería Informática y de Sistemas. Para ello, se establece el siguiente protocolo de uso de infraestructura acordado con la Dirección de la Escuela y la Unidad de Infraestructura de la UNSAAC:

\subsection*{Uso de Espacios Técnicos (Cuartos de Telecomunicaciones)}

Se formaliza la autorización para la habilitación exclusiva de los espacios identificados en los planos de diseño como MDF (Main Distribution Frame) del Nivel 2 e IDF (Intermediate Distribution Frame) de los Niveles 3 y 4.

\begin{itemize}
    \item \textbf{Derecho de Uso:} El contratista está facultado para instalar gabinetes, anclajes antisísmicos y organizadores verticales en las áreas designadas para el MDF del Nivel 2 y los IDF-03 e IDF-04.
    \item \textbf{Restricción de Acceso:} Durante la fase de implementación, estos espacios quedarán bajo custodia del contratista. Una vez entregada la obra, se entregarán las llaves y códigos de acceso a la Jefatura de Laboratorios o al Centro de Cómputo de la Facultad.
\end{itemize}

\subsection*{Intervención en Rutas Compartidas (Montantes y Pasillos)}

Se autoriza la intervención en las zonas comunes para el tendido del cableado estructurado, bajo las siguientes condiciones de convivencia con otros servicios:

\begin{enumerate}
    \item \textbf{Montantes Verticales (Shafts):} Permiso para utilizar los ductos verticales existentes que conectan los pisos 1 al 4.
    \begin{itemize}
        \item \textbf{Condición:} Se respetará rigurosamente el espacio ocupado por las instalaciones eléctricas existentes, instalando separadores físicos o tubería conduit independiente para garantizar la inmunidad al ruido electromagnético.
    \end{itemize}
    \item \textbf{Falsos Techos y Pasadizos:} Autorización para el desmontaje temporal de baldosas en los pasillos de circulación para la fijación de soportes y canaletas. El contratista se compromete a restituir cualquier baldosa dañada durante el proceso.
\end{enumerate}

\subsection*{Protocolo de Intervención en Áreas Académicas (Aulas y Laboratorios)}

Dado que los trabajos se realizarán en ambientes de enseñanza en los \totalPisos{} niveles (Aulas 101\,--\,107 del Nivel 1, Aulas 201\,--\,208 del Nivel 2, Laboratorios 301\,--\,312 del Nivel 3, Cubículos y Auditorio del Nivel 4), el permiso de uso está condicionado al cronograma académico:

\begin{itemize}
    \item \textbf{Ventanas de Trabajo:} Las perforaciones ruidosas y el tendido de cable en aulas se realizarán prioritariamente en horarios de baja afluencia (tardes, fines de semana o periodos vacacionales) para no interrumpir el dictado de clases.
    \item \textbf{Liberación de Ambientes:} El contratista deberá solicitar la disponibilidad del aula con 48 horas de anticipación y entregarla limpia y operativa antes del inicio de la primera clase del día siguiente.
\end{itemize}

\section{Documentos de propiedad o administración}

Para garantizar la seguridad jurídica del proyecto y la correcta asignación de responsabilidades sobre la infraestructura intervenida, se integrarán al expediente los siguientes documentos administrativos que acreditan la titularidad y la gestión del inmueble:

\subsection*{Acreditación de Titularidad y Saneamiento Físico-Legal}

Se hace constar que el área de intervención (Pabellón de la Escuela Profesional de Ingeniería Informática y de Sistemas) se encuentra dentro del predio matriz de la Ciudad Universitaria de Perayoc, propiedad inscrita y saneada a favor de la Universidad Nacional de San Antonio Abad del Cusco (UNSAAC).

\begin{itemize}
    \item \textbf{Base Legal:} La intervención se ampara en los documentos de gestión patrimonial de la Universidad, lo que faculta a la Oficina de Infraestructura a autorizar mejoras y mantenimientos dentro de sus instalaciones sin requerir permisos de terceros propietarios.
\end{itemize}

\subsection*{Actos Administrativos de Aprobación (Resoluciones)}

La ejecución financiera y técnica del proyecto deberá estar respaldada por los actos administrativos correspondientes emitidos por el Consejo de Facultad o la autoridad competente:

\begin{enumerate}
    \item \textbf{Resolución de Aprobación de Expediente Técnico:} Documento oficial que aprueba el presente diseño, presupuesto y cronograma, autorizando el gasto o la inversión.
    \item \textbf{Resolución de Designación de Supervisión:} Documento que nombra formalmente al Inspector o Supervisor de obra por parte de la Facultad, quien tendrá la autoridad para firmar las valorizaciones y el acta de recepción final.
\end{enumerate}

\subsection*{Acta de Entrega de Terreno}

Antes del ``Día 0'' del cronograma, se requiere la suscripción del Acta de Entrega de Terreno.

\begin{itemize}
    \item \textbf{Propósito:} Mediante este documento legal, la administración del edificio (Decanato/Dirección de Escuela) entrega temporalmente la posesión de los ambientes (aulas, laboratorios y cuartos técnicos) al Ejecutor.
    \item \textbf{Efecto:} A partir de la firma de esta acta, la responsabilidad por la seguridad, limpieza y custodia de los bienes dentro de las áreas de trabajo recae exclusivamente sobre el Ejecutor hasta la recepción final de la obra.
\end{itemize}

\section{Contratos y acuerdos}

Para asegurar la sostenibilidad del proyecto y formalizar la relación entre el Ejecutor y la Facultad de Ingeniería Informática y de Sistemas, se establecen los siguientes instrumentos contractuales y entregables técnicos:

\subsection*{Autorizaciones de obra interna en la edificación}

Se suscribe el Acuerdo de Intervención en Planta Interna, mediante el cual la Universidad autoriza al Ejecutor a realizar modificaciones no estructurales en los cuatro niveles del pabellón.

\begin{itemize}
    \item \textbf{Alcance del Acuerdo:} Permite la fijación de canaletas adosadas, perforación de muros divisorios y desmontaje temporal de falsos techos.
    \item \textbf{Cláusula de Responsabilidad Civil:} El Ejecutor asume plena responsabilidad contractual por cualquier daño incidental causado a la infraestructura (ej. rotura de vidrios, daño a mobiliario o afectación de acabados), comprometiéndose a la restitución inmediata del bien a su estado original sin costo para la Universidad.
\end{itemize}

\subsection*{Permisos de uso de infraestructura existente}

Se formaliza el Contrato de Co-ubicación y Uso de Vías, el cual regula el aprovechamiento de los recursos físicos actuales del edificio.

\begin{itemize}
    \item \textbf{Ductos y Montantes:} Se concede el derecho de paso por los ductos verticales (shafts) existentes, bajo la condición de que el nuevo cableado Cat 6A se instale debidamente peinado y etiquetado, respetando el factor de llenado máximo del 40\% según norma TIA-569-D.
    \item \textbf{Energía y Tierras:} Se autoriza la conexión a los tableros eléctricos designados y a la malla de tierra del edificio para la protección de los gabinetes, previo visto bueno de la Oficina de Mantenimiento Eléctrico de la UNSAAC.
\end{itemize}

\subsection*{Manual de operación y mantenimiento de la red instalada}

Como parte fundamental del cierre contractual, el Ejecutor entregará el Manual Técnico de la Red, un documento guía diseñado para el administrador de la red (Network Admin) de la Facultad. Este manual contendrá:

\begin{enumerate}
    \item \textbf{Documentación de Ingeniería (As-Built):}
    \begin{itemize}
        \item Planos finales en AutoCAD con la ubicación exacta de los \textbf{\totalPuntos{}} puntos y la ruta de canalización real en los \totalPisos{} niveles.
        \item Esquema lógico de distribución de los Racks (Front-View) indicando qué panel atiende a qué laboratorio.
    \end{itemize}
    \item \textbf{Guía de Mantenimiento Preventivo:}
    \begin{itemize}
        \item Procedimientos para la limpieza correcta de conectores y patch cords.
        \item Recomendaciones para la gestión de cambios (MACs - Moves, Adds, Changes) para mantener la garantía vigente.
        \item Instrucciones para la lectura e interpretación del etiquetado TIA-606-C implementado.
    \end{itemize}
    \item \textbf{Protocolo de Solución de Problemas (Troubleshooting):}
    \begin{itemize}
        \item Árbol de decisión para diagnosticar fallas comunes (ej. ``Usuario sin red en Lab 301'').
        \item Datos de contacto del soporte técnico del instalador para atención de garantías.
    \end{itemize}
\end{enumerate}
